\chapter{समानुपात, चाल और औसत}\label{ch:g8:speed}

अब \omterm{def:g7:prop:table}{समानुपात} (\cref{ch:g7:prop}) की मुलाक़ात असली दुनिया से होती है:
\omterm{def:g8:speed:def}{चाल}, बहाव-दर, मात्रक बदलना --- यानी ऐसी राशियाँ जो दो और राशियों के
\emph{\omterm{def:g3:division:remainder}{भागफल}} हैं। अध्याय भारित औसतों पर ख़त्म होता है, जहाँ
\omterm{def:g7:prop:table}{समानुपात} वाली सोच एक जानी-पहचानी ग़लती से बचा लेती है।

\section{औसत चाल}

\begin{definition}[औसत चाल]\label{def:g8:speed:def}
किसी सफ़र की \emph{औसत चाल}\index{चाल} यह \omterm{def:g3:division:remainder}{भागफल} है
\[
v = \frac{d}{t}
\qquad
\left(\text{दूरी बटा अवधि}\right),
\]
और उसका मात्रक किलोमीटर प्रति घंटा (km/h), मीटर प्रति सेकंड (m/s),
\dots\ होता है। यही सूत्र उलटकर भी चलता है: $d = v \times t$ और
$t = \frac{d}{v}$।
\end{definition}

\begin{example}\label{ex:g8:speed:basic}
एक रेलगाड़ी $1$ h $30$ min में $270$ km तय करती है। पहले समय बदलो:
$1$ h $30$ min $= 1.5$ h। फिर
\[
v = \frac{270}{1.5} = 180 \text{ km/h}.
\]
$90$ km/h से $2$ h $20$ min $= \frac73$ h में तय दूरी:
$d = 90 \times \frac73 = 210$ km। और $14$ km/h से $35$ km में लगा
समय: $t = \frac{35}{14} = 2.5$ h $= 2$ h $30$ min।
\end{example}

\begin{remark}[मिनट दशमलव नहीं होते]
$1$ h $30$ min, $1.5$ h है, पर $1$ h $20$ min, $1.2$ h \emph{नहीं}
है: वह $1 + \frac{20}{60} = \frac43 \approx 1.33$ h है। भाग देने से
पहले मिनटों को हमेशा घंटे की \omterm{def:g6:fractions:def}{भिन्न} में बदलो ($60$ min $= 1$ h)।
\end{remark}

\begin{omfigure}
\begin{tikzpicture}
\begin{axis}[omaxis,
  width=8.4cm, height=5.6cm,
  xmin=0, xmax=3.4, ymin=0, ymax=270,
  xlabel={समय (h)}, ylabel={दूरी (km)},
  xtick={1,2,3}, ytick={80,160,240}]
  \addplot[omDef, thick, domain=0:3.2] {80*x};
  \addplot[omThm, thick, domain=0:3.2] {50*x};
  \node[omDef, font=\small, rotate=36] at (axis cs:2.35,215)
    {$80$ km/h};
  \node[omThm, font=\small, rotate=25] at (axis cs:2.6,155)
    {$50$ km/h};
\end{axis}
\end{tikzpicture}

{\small एक जैसी \omterm{def:g8:speed:def}{चाल} पर दूरी समय के \omterm{def:g6:propdata:def}{समानुपाती} होती है: मूलबिंदु से
जाती एक सीधी रेखा, जिसकी ढलान \emph{ही} \omterm{def:g8:speed:def}{चाल} है।}
\end{omfigure}

\begin{method}[चाल के मात्रक बदलना]\label{met:g8:speed:convert}
km/h को m/s में बदलने के लिए: $1$ km $= 1000$ m और $1$ h $= 3600$ s,
इसलिए
\[
v \text{ km/h} = \frac{v \times 1000}{3600} \text{ m/s}
= \frac{v}{3.6} \text{ m/s}.
\]
m/s को km/h में बदलने के लिए $3.6$ से गुणा करो।
\end{method}

\begin{example}\label{ex:g8:speed:convert}
$90$ km/h $= \frac{90}{3.6} = 25$ m/s। $100$ m $10$ s में दौड़ने
वाला धावक $10$ m/s $= 36$ km/h से चलता है।
\end{example}

\section{भागफल वाली राशियाँ}

\begin{example}[बहाव, घनत्व, प्रति किलो दाम]\label{ex:g8:speed:quotients}
\omterm{def:g8:speed:def}{चाल} के कई चचेरे भाई हैं, और सबका बरताव एक जैसा है:
\begin{itemize}
  \item एक नल $4$ min में $48$ L भरता है: बहाव-दर
        $\frac{48}{4} = 12$ L/min, इसलिए $150$ L की एक टंकी भरने
        में $\frac{150}{12} = 12.5$ min लगते हैं;
  \item $0.6$ kg पनीर $9$ यूरो का है: दाम
        $\frac{9}{0.6} = 15$ यूरो प्रति kg;
  \item एक गाड़ी $90$ km में $6.3$ L ख़र्च करती है: खपत
        $\frac{6.3}{90} \times 100 = 7$ L प्रति $100$ km।
\end{itemize}
\omterm{def:g3:division:remainder}{भागफल} पहचानो, और तीन तरफ़ा सूत्र ($q = \frac ab$, $a = qb$,
$b = \frac aq$) बाक़ी काम कर देता है।
\end{example}

\section{भारित औसत}

\begin{definition}[भारित औसत]\label{def:g8:speed:average}
जब मानों $v_1, v_2, \dots$ के साथ गिनतियाँ (या भार)
$n_1, n_2, \dots$ आती हों, तो उनका \emph{भारित औसत}\index{भारित माध्य}
यह है
\[
\bar v = \frac{n_1 v_1 + n_2 v_2 + \dots}{n_1 + n_2 + \dots} :
\]
यानी मानों का कुल, भारों के कुल से भाग दिया हुआ।
\end{definition}

\begin{example}\label{ex:g8:speed:average}
एक परीक्षा दो समूहों ने दी: समूह 1 में $12$ विद्यार्थी, औसत $11$;
समूह 2 में $18$ विद्यार्थी, औसत $14$। पूरी कक्षा का औसत:
\[
\bar v = \frac{12 \times 11 + 18 \times 14}{12 + 18}
= \frac{132 + 252}{30} = \frac{384}{30} = 12.8 .
\]
यह $\frac{11 + 14}{2} = 12.5$ \emph{नहीं} है: बड़ा समूह औसत को अपनी
तरफ़ खींच लेता है --- भार मायने रखते हैं।
\end{example}

\begin{example}[दो चरणों की औसत चाल]\label{ex:g8:speed:twolegs}
एक साइकिल सवार $30$ km, $30$ km/h से चलता है, फिर $30$ km,
$15$ km/h से। पूरे सफ़र की \omterm{def:g8:speed:def}{औसत चाल}
$\frac{30 + 15}{2} = 22.5$ km/h \emph{नहीं} है। परिभाषा से निकालो:
\begin{enumerate}
  \item समय: $\frac{30}{30} = 1$ h, फिर $\frac{30}{15} = 2$ h;
        कुल $3$ h;
  \item कुल दूरी $60$ km, इसलिए
        $v = \frac{60}{3} = 20$ km/h।
\end{enumerate}
धीमा चरण ज़्यादा देर चलता है, इसलिए उसका भार ज़्यादा है --- चालों का
औसत हमेशा \emph{समय} के भार से निकालना पड़ता है।
\end{example}

\section{अभ्यास}

\begin{exercise}[$\star$]\label{exo:g8:speed:1}
\omterm{def:g8:speed:def}{औसत चाल} निकालो: $2$ h में $150$ km; $45$ min में $27$ km;
$12.5$ s में $100$ m।
\end{exercise}

\begin{exercise}[$\star$]\label{exo:g8:speed:2}
दूरी निकालो: $85$ km/h से $2$ h; $90$ km/h से $40$ min। समय
निकालो: $90$ km/h से $315$ km (h और min में)।
\end{exercise}

\begin{exercise}[$\star$]\label{exo:g8:speed:3}
बदलो: $72$ km/h को m/s में; $5$ m/s को km/h में; $108$ km/h को m/s में।
\end{exercise}

\begin{exercise}[$\star$]\label{exo:g8:speed:4}
एक नल $15$ L/min देता है। $600$ L की टंकी भरने में कितना समय लगेगा?
और $2$ h $30$ में कितने लीटर?
\end{exercise}

\begin{exercise}[$\star$]\label{exo:g8:speed:5}
कौन-सा सौदा बेहतर है: $1.2$ kg सेब $3.30$ यूरो में, या $0.8$ kg
$2.16$ यूरो में? प्रति किलोग्राम दाम तुलो।
\end{exercise}

\begin{exercise}[$\star$]\label{exo:g8:speed:6}
$25$ विद्यार्थियों की एक कक्षा का किसी परीक्षा में औसत $12.4$ है;
$35$ विद्यार्थियों की दूसरी कक्षा का औसत $10$। दोनों कक्षाओं का
मिलाकर औसत निकालो।
\end{exercise}

\begin{exercise}[$\star$]\label{exo:g8:speed:7}
आवाज़ क़रीब $340$ m/s चलती है। तुम बिजली की कौंध देखते हो और $6$
सेकंड बाद गरज सुनते हो। बिजली कितनी दूर गिरी (नज़दीकी $100$ m तक)?
\end{exercise}

\begin{exercise}[$\star\star$]\label{exo:g8:speed:8}
एक पैदल यात्री $2$ h, $5$ km/h से चलता है, $30$ min सुस्ताता है,
फिर $1$ h $30$, $4$ km/h से चलता है। कुल दूरी और \emph{सुस्ताने
समेत} \omterm{def:g8:speed:def}{औसत चाल} निकालो (कुल दूरी बटा कुल बीता समय)।
\end{exercise}

\begin{exercise}[$\star\star$]\label{exo:g8:speed:9}
गुणांकों वाले अंक: किसी विद्यार्थी को $15$ (गुणांक $3$), $9$
(गुणांक $2$) और $12$ (गुणांक $1$) मिले। \omterm{def:g8:speed:average}{भारित औसत} निकालो। सादा
(बिना भार का) औसत क्या होता, और दोनों अलग क्यों हैं?
\end{exercise}

\begin{exercise}[$\star\star$]\label{exo:g8:speed:10}
एक गाड़ी $120$ km, $80$ km/h से चलती है और फिर $120$ km, $120$ km/h से।
\begin{enumerate}
  \item हर चरण की अवधि निकालो, फिर पूरे सफ़र की \omterm{def:g8:speed:def}{औसत चाल}।
  \item समझाओ कि उत्तर दोनों चालों के बीच वाले $100$ km/h से कम
        क्यों है।
\end{enumerate}
\end{exercise}

\begin{exercise}[$\star\star\star$]\label{exo:g8:speed:11}
दो गाँव $18$ km दूर हैं। अन्ना पहले गाँव से $10{:}00$ बजे निकलकर
दूसरे की ओर $5$ km/h से चलती है; बोरिस उसी समय दूसरे गाँव से पहले की
ओर $4$ km/h से चलता है। वे किस समय मिलेंगे, और अन्ना के गाँव से
कितनी दूर? (दोनों मिलकर फ़ासला $5 + 4$ km/h की दर से घटाते हैं।)
\end{exercise}

\section{समस्या: हरात्मक माध्य, यानी वापसी का सफ़र औसत क्यों बिगाड़
देता है}

\begin{problem}[{सप्ताहांत समस्या --- बराबर दूरियों पर \omterm{def:g8:speed:def}{औसत चाल}
$\frac{2 v_1 v_2}{v_1 + v_2}$ होती है, और वह दोनों चालों के बीच
वाले मान से कभी आगे नहीं निकलती}]
\label{pb:g8:speed:1}
\Cref{ex:g8:speed:twolegs} और \cref{exo:g8:speed:10} दोनों एक ही
हैरानी पर ख़त्म हुए थे: एक \omterm{def:g8:speed:def}{चाल} से जाओ, दूसरी से लौटो, और \omterm{def:g8:speed:def}{औसत चाल}
दोनों के बीच वाले मान से \emph{नीचे} जा गिरती है। यह समस्या उस
हैरानी के पीछे का ठीक-ठीक सूत्र ढूँढ़ती है --- एक नई तरह का औसत,
\emph{हरात्मक माध्य}\index{हरात्मक माध्य} --- सिद्ध करती है कि यह
हैरानी दरअसल एक प्रमेय है, और उसे उसके तार्किक सिरे तक ले जाती है:
एक ऐसी भरपाई जो दुनिया की कोई भी \omterm{def:g8:speed:def}{चाल} कर ही नहीं सकती।

\medskip
\noindent\textbf{भाग I --- आने-जाने के सफ़र का सूत्र।}
किसी सफ़र में दूरी $d$ \omterm{def:g8:speed:def}{चाल} $v_1$ से तय होती है, फिर उतनी ही दूरी $d$
\omterm{def:g8:speed:def}{चाल} $v_2$ से (सारी संख्याएँ धन हैं)।
\begin{enumerate}
  \item दोनों अवधियाँ, फिर कुल अवधि, $d$, $v_1$, $v_2$ वाली
        भिन्नों के रूप में लिखो।
  \item पूरे सफ़र की \omterm{def:g8:speed:def}{औसत चाल}
        $\bar v = \dfrac{2d}{\ \dfrac{d}{v_1} + \dfrac{d}{v_2}\ }$
        है --- यानी एक दुमंज़िला \omterm{def:g6:fractions:def}{भिन्न}
        (\cref{ex:g8:fractions:complex})। उसे सरल करो और सिद्ध
        करो:
        \[
        \bar v = \frac{2\, v_1 v_2}{v_1 + v_2} .
        \]
  \item ऊपर बताई गई दोनों गणनाओं पर यह सूत्र जाँचो: $v_1 = 30$,
        $v_2 = 15$ km/h (\cref{ex:g8:speed:twolegs}), और
        $v_1 = 80$, $v_2 = 120$ km/h (\cref{exo:g8:speed:10})।
  \item सूत्र में से दूरी $d$ ग़ायब हो गई है। साइकिल सवार के लिए
        इसका \omterm{def:g5:solids:def}{ठोस} मतलब क्या है?
  \item अब मान लो सफ़र में हर \omterm{def:g8:speed:def}{चाल} पर एक ही \emph{समय} $T$ बीतता है।
        दिखाओ कि तब \omterm{def:g8:speed:def}{औसत चाल} $\frac{v_1 + v_2}{2}$ होती है, यानी
        सादा बीच वाला मान। \cref{def:g8:speed:average} की भाषा
        में: चालों का औसत हमेशा किस राशि के भार से निकालना पड़ता
        है, और दोनों स्थितियों में बदलता क्या है?
\end{enumerate}

\medskip
\noindent\textbf{भाग II --- हरात्मक बनाम \omterm{def:g6:lines:perp}{समांतर}।}
दो धन संख्याओं $v_1$, $v_2$ के लिए मान लो
\[
H = \frac{2\, v_1 v_2}{v_1 + v_2}
\quad\text{(उनका \emph{हरात्मक माध्य}),}
\qquad
A = \frac{v_1 + v_2}{2}
\quad\text{(उनका \emph{समांतर माध्य})।}
\]
\begin{enumerate}[resume]
  \item जोड़ों $(30, 15)$ और $(80, 120)$ के लिए $H$ तथा $A$
        निकालो। हर बार दोनों माध्यों में से कौन जीतता है?
  \item $A - H$ को उभयनिष्ठ हर $2(v_1 + v_2)$ पर लाओ, \omterm{def:g4:fractions:def}{अंश} को
        उल्लेखनीय सर्वसमिकाओं (\cref{pb:g8:equations:1}) से
        फैलाओ, और सिद्ध करो:
        \[
        A - H = \frac{(v_1 - v_2)^2}{2\,(v_1 + v_2)} .
        \]
  \item इससे सारी हैरानियों के पीछे बैठी प्रमेय निकालो: धन चालों
        के लिए हमेशा $H \leq A$, और बराबरी ठीक तब जब
        $v_1 = v_2$। \omterm{def:g4:fractions:def}{अंश} में बैठा एक वर्ग यह मामला क्यों तय कर
        देता है?
  \item यह सुंदर सर्वसमिका सिद्ध करो: $H \times A = v_1 \times v_2$
        --- दोनों माध्य गुणा होकर वापस असली \omterm{def:g2:mult:def}{गुणनफल} बना देते हैं।
  \item एक गाड़ी $60$ km, $40$ km/h से और $60$ km, $60$ km/h से तय
        करती है। सूत्र से \omterm{def:g8:speed:def}{औसत चाल} बताओ, फिर लंबे रास्ते से ---
        समय और कुल दूरी निकालकर --- उसकी पुष्टि करो।
\end{enumerate}

\medskip
\noindent\textbf{भाग III --- वह भरपाई जो हो ही नहीं सकती।}
\begin{enumerate}[resume]
  \item एक साइकिल सवार $30$ km की दौड़ में $30$ km/h का औसत रखने की
        उम्मीद से उतरती है। दौड़ ज़्यादा से ज़्यादा कितने समय की हो
        सकती है? वह पहले $15$ km, $15$ km/h से चलती है: उस समय के
        बजट में कितना बचा?
  \item दिखाओ कि उम्मीद वाला औसत अब सचमुच नामुमकिन हो चुका है:
        दूसरा \omterm{def:g3:division:half}{आधा} हिस्सा $90$ km/h से चलने पर उसकी \omterm{def:g8:speed:def}{औसत चाल}
        निकालो, और समझाओ कि कितनी भी बड़ी \omterm{def:g8:speed:def}{चाल} क्यों उसे $30$ km/h
        तक नहीं पहुँचा सकती।
  \item वह अपना लक्ष्य घटाकर $20$ km/h कर देती है। बाक़ी $15$ km
        पर कौन-सी \omterm{def:g8:speed:def}{चाल} यह लक्ष्य पूरा करेगी? अपना उत्तर \omterm{pb:g8:speed:1}{हरात्मक
        माध्य} वाले सूत्र से जाँचो।
  \item अकेला ठंडे पानी का नल एक टब $20$ min में भरता है; अकेला
        गरम पानी का नल $30$ min में। दोनों नल मिलकर एक मिनट में
        टब का कितना भाग भरते हैं, और पूरा टब भरने में कितना समय
        लेते हैं? अपने उत्तर की तुलना $H(20, 30)$ से करो ---
        तुम्हें क्या दिखता है?
  \item एक हवाई जहाज़ एक ही रास्ते पर जाता और लौटता है: हवा के साथ
        $900$ km/h, हवा के ख़िलाफ़ $700$ km/h। आने-जाने की \omterm{def:g8:speed:def}{औसत चाल}
        निकालो, और एक वाक्य में समझाओ कि हवा चाहे जैसी चले, वह
        आने-जाने के सफ़र को हमेशा \emph{लंबा} ही क्यों करती है।
\end{enumerate}
\end{problem}
